%==============================================================
%  Félix Jesús Nilo Bustios — Currículum Vitae (formato Harvard)
%  Compilar con: pdflatex Felix_Nilo_Bustios_CV_ES.tex
%  (o pegar en https://overleaf.com — no se requieren paquetes extra)
%==============================================================
\documentclass[11pt,a4paper]{article}

\usepackage[T1]{fontenc}
\usepackage[utf8]{inputenc}
\usepackage[spanish]{babel}
\usepackage[margin=1.6cm]{geometry}
\usepackage{enumitem}
\usepackage{titlesec}
\usepackage{xcolor}
\usepackage[hidelinks]{hyperref}

\definecolor{accent}{HTML}{1F3A5F}

% --- encabezados de sección: versalitas, negrita, línea -------
\titleformat{\section}
  {\normalsize\bfseries\scshape\color{accent}}{}{0em}{}[\vspace{1pt}\hrule height 0.8pt]
\titlespacing{\section}{0pt}{8pt}{5pt}

% --- configuración de listas ----------------------------------
\setlist[itemize]{leftmargin=1.25em, itemsep=1.5pt, topsep=2pt, parsep=0pt, label=\textbullet}

\setlength{\parindent}{0pt}
\pagestyle{empty}

% --- macro de entrada: {Organización}{Ubicación}{Cargo}{Fechas} ------
\newcommand{\entry}[4]{%
  \textbf{#1}\hfill\textbf{#2}\par
  \textit{#3}\hfill\textit{#4}\par\vspace{2pt}%
}

\begin{document}

%================= ENCABEZADO =================================
\begin{center}
  {\LARGE\bfseries\color{accent} Félix Jesús Nilo Bustios}\\[4pt]
  {\small
    \href{mailto:felixjesusnilob@gmail.com}{felixjesusnilob@gmail.com}\\
    \href{https://www.linkedin.com/in/felix-nilo}{LinkedIn} \ \textbullet\
    \href{https://github.com/felixnilo}{GitHub}
  }
\end{center}

%================= RESUMEN ====================================
\section{Resumen}

Arquitecto AWS y líder de ingeniería con más de 6 años construyendo sistemas cloud orientados
a eventos y equipos de ingeniería de alto rendimiento. En OptConnect, lidero la migración de
integraciones de operadoras desde un monolito Perl legado hacia una arquitectura serverless
moderna conectando Verizon, Vodafone, Webbing y Teal a través de un pipeline multi-etapa
orientado a eventos (Lambda $\to$ DynamoDB Streams $\to$ MySQL $\to$ Rails UI) que gestiona
tráfico real de producción para miles de dispositivos IoT.

Impacto medible: \textbf{$-$40\%} en errores de producción, \textbf{$\sim$50\%} de reducción en
costos AWS, \textbf{85\%} de cumplimiento de OKRs, y el crecimiento profesional de ingenieros
en \textbf{+9 niveles de seniority combinados}.

Manos en la masa con Python, Serverless Framework, arquitectura orientada a eventos y Ruby on Rails.
Opero en la intersección del diseño de sistemas distribuidos, liderazgo de equipos y resultados de negocio.

%================= EXPERIENCIA ================================
\section{Experiencia}

\entry{OptConnect}{Santiago, Chile}{Desarrollador Python AWS}{Mayo 2025 -- Presente}
\begin{itemize}
  \item Lideré la migración de \textbf{+10 flujos críticos de operadoras} desde un monolito Perl
        legado hacia una arquitectura serverless orientada a eventos (Lambda, API Gateway,
        EventBridge, SQS, SNS, DynamoDB Streams, RDS, VPC/Endpoints), conectando
        \textbf{Verizon, Vodafone, Webbing y Teal} vía REST y SOAP en Serverless Framework.
  \item Contribuí a una \textbf{reducción de $\sim$50\% en la factura AWS} mediante la migración
        serverless y optimizaciones de costos como la recuperación cacheada de Secrets Manager
        (AWS Powertools).
  \item Re-arquitecté el pipeline de callbacks de operadoras como un sistema distribuido e
        idempotente: los callbacks de Verizon llegan a API Gateway $\to$ Lambda $\to$ DynamoDB
        (put condicional, dedup) $\to$ SNS notifica a Lola; una rama paralela de DynamoDB Streams
        fluye por un hub de agregación $\to$ MySQL como fuente de verdad durable --- cubriendo
        rutas de éxito, fallo y reintento, con retención basada en TTL, validado end-to-end
        contra tráfico real de producción.
  \item Diseñé y entregué \textbf{NotificationService}, el punto de entrada único para callbacks
        de Verizon ($\sim$\textbf{1.000/hora} en producción): API Gateway + Lambda reenvía cada
        callback crudo vía HTTPS firmado con SigV4 a sus suscriptores por tipo de callback
        (VerizonIntegration, la UI Rails, el pipeline de datos). Reemplazar el fan-out de SNS
        eliminó el límite de 256\,KB que obligaba a dividir el \textbf{96\%} de los callbacks de
        uso; la entrega todo-o-nada devuelve los reintentos a la operadora, con trazabilidad y
        alertas por callback.
  \item Contribuí al \textbf{pipeline de datos GCS} de la empresa (Extract $\to$ Transform $\to$
        Load: S3, Kinesis, PostgreSQL), enrutando callbacks de operadoras hacia él y diseñando
        \textbf{gcs-commander}, el servicio que enviará comandos de ciclo de vida de dispositivos
        (suspender, restaurar, activar, terminar) a las operadoras y reemplazará la integración
        Verizon: semántica de envío único contra acciones duplicadas de la operadora, entrega
        asíncrona en cola, y resultados cerrados por callbacks de la operadora o un deadline.
        Implementé su contrato y el comando de referencia para Verizon.
  \item Ejecuté análisis de uso y calidad de datos en SQL sobre la base de datos MySQL de
        producción entre operadoras, para facturación y operaciones: tracé incidentes de
        integridad hasta su causa raíz en el pipeline --- un mes de uso reescrito $\sim$6\%
        por debajo en \textbf{$\sim$91k líneas} por callbacks re-entregados, planes de tarifa
        desactualizados en \textbf{$\sim$550k dispositivos} --- y audité tendencias de uso por
        cuenta y SIMs activas ociosas por cliente.
  \item Construí automatizaciones Python en Odoo (ERP) dando visibilidad al negocio sobre
        \textbf{$\sim$50.000 transacciones mensuales}, integrando datos de pago para contabilidad
        y contactos.
  \item Diseñé un bot interno de IA en Slack sobre Amazon Bedrock (Nova Lite) + Lambda/API
        Gateway/Secrets Manager: sanitización de entrada, kill switch, contexto multi-turno,
        logging estructurado y dashboard en CloudWatch para volumen de consultas, latencia y errores.
  \item Automaticé infraestructura (CloudFormation + Make), reduciendo el tiempo de configuración
        para nuevos desarrolladores para que el equipo se enfoque en valor y no en configuración.
  \item Construí CI/CD para sistemas serverless y pipelines de invocación externa; plantillé el
        Serverless Framework con auto-PRs de GitHub Actions sincronizados entre repositorios.
  \item Construí y mantuve dashboards de Grafana en múltiples entornos (fuente de datos
        CloudWatch, Assume Role), mejorando la visibilidad operacional de los sistemas de
        integración de operadoras.
  \item Responsable del diseño de sistemas end-to-end: propuse y evalué arquitecturas (Step
        Functions, SQS fan-out, VPC Endpoints), produje especificaciones técnicas y coordiné
        la entrega entre equipos onshore y offshore.
\end{itemize}

\entry{Innovación y Tecnología Aplicada SpA}{La Serena, Chile}{Consultor de Producto e Ingeniería}{Sep 2024 -- Presente}
\begin{itemize}
  \item Realicé descubrimiento de producto (1:1 con clientes clave de la industria) y
        \textbf{digitalicé procesos core de empresas regionales}, volcando los insights en
        mockups de Figma y roadmap de producto.
  \item Conduje \textbf{evaluaciones de seguridad OSINT descubriendo vulnerabilidades en 50
        empresas regionales}.
  \item Entregué varias apps web y móviles en producción con React Native (Google Play \& App
        Store), React con Vite y backend Django 5 REST con pipeline CI/CD en GitHub Actions.
  \item Orquesté entornos de producción, staging y desarrollo en stacks serverless (AWS Lambda,
        API Gateway, S3, DynamoDB, RDS y CloudFormation) y aplicaciones dockerizadas desplegadas
        en AWS con balanceo de carga para alta disponibilidad.
  \item Integré Google/Apple Auth, Google Ads, Pagos Google/iOS, MediaPipe SDK, Posthog y Sentry;
        tomé decisiones con analytics de MS Clarity y pruebas A/B.
  \item Gestioné el equipo de desarrollo, proceso de onboarding, backlog y planes de carrera.
  \item Configuré Cloudflare y CloudFront para entrega de contenido y seguridad, garantizando
        acceso rápido y seguro a las aplicaciones.
  \item Definí especificaciones para soluciones digitales aplicadas en \textbf{agricultura,
        fitness, minería entre otros} sectores, incluyendo apps móviles y web para
        \textbf{Gestión de Cultivos, Maximización de Conversiones y Cumplimiento Ambiental}
        con datos UAV.
  \item Alcancé y desarrollé un \href{https://amissionhub.com}{\textbf{Sistema de Gestión de
        Cultivos}} con datos UAV, incluyendo monitoreo de salud de cultivos, predicción de
        rendimiento y técnicas de agricultura de precisión.
  \item Desarrollé una \href{https://dev.blou.club/}{\textbf{App de Fitness}} con planes de
        entrenamiento personalizados, seguimiento de progreso y funciones sociales para mejorar
        el engagement y motivación del usuario.
  \item Desarrollé una \href{https://chemical2.com/}{\textbf{App de Cumplimiento Ambiental}}
        para empresas mineras, con monitoreo en tiempo real de parámetros ambientales, reportes
        y seguimiento de cumplimiento normativo.
  \item Implementé una \href{https://dev.blou.club/}{\textbf{App de Maximización de
        Conversiones}} para empresas, con pruebas A/B, analytics y herramientas de optimización
        para mejorar tasas de conversión e impulsar el crecimiento del negocio.
\end{itemize}

\entry{Universidad Católica del Norte / Universidad Central}{Coquimbo / La Serena, Chile}{Docente Ingeniería Civil en Computación e Informática}{Jun 2022 -- Dic 2025}
\begin{itemize}
  \item Diseñé syllabi y facilité \textbf{5 cursos de pregrado y 1 de postgrado} en dos universidades, abarcando Ingeniería de Software, Proyectos Aplicados y fundamentos de sistemas.
  \item \textbf{Universidad Católica del Norte}: facilité \textit{Introducción a la Ingeniería}, fomentando la identidad profesional y el pensamiento ingenieril en estudiantes de primer año.
  \item \textbf{Universidad Central}: facilité \textit{Arquitectura de Software}, \textit{Laboratorio de Complejidad de Algoritmos}, \textit{Fundamentos de Lenguajes de Programación} y \textit{Tópicos Avanzados de Datos} (pregrado); y \textit{Virtualización} (postgrado) --- articulando teoría con práctica alineada a la industria.
\end{itemize}


\entry{BUK SpA}{Santiago, Chile}{Engineering Manager}{Feb 2023 -- Sep 2024}
\begin{itemize}
  \item Hice crecer al equipo \textbf{6 niveles de seniority combinados} mediante mentoría,
        feedback continuo y evaluaciones de desempeño.
  \item Alcancé el \textbf{85\% de los OKRs del equipo} y reduje los errores de producción en
        \textbf{40\%} mediante una estrategia de mejora de calidad.
  \item Entregué \textbf{10 funcionalidades de alto impacto} para el módulo de Gestión del
        Desempeño y el Marketplace de Soluciones e Integraciones usando Ruby on Rails.
  \item Responsable del reporte a stakeholders sobre desempeño del equipo, priorización del
        roadmap y operaciones críticas.
\end{itemize}

\entry{Arkho Technology SpA}{Santiago, Chile}{Delivery Manager \& Team Lead}{Feb 2022 -- Feb 2023}
\begin{itemize}
  \item Gestioné \textbf{4 equipos de operaciones en paralelo}, haciendo crecer el seniority
        \textbf{3 niveles}.
  \item Entregué \textbf{7+ propuestas técnico-comerciales cloud} y conduje \textbf{10 procesos
        de onboarding} en squads, proyectos y servicios gestionados.
  \item Apliqué AWS Step Functions, S3, Fargate, Lightsail, WAF, Route\,53, Lambda, EC2, API
        Gateway, Glue y QuickSight; lideré planificación con stakeholders en métricas de calidad
        y productividad.
\end{itemize}

\entry{ES Ético y Sustentable Ltda.}{La Serena, Chile}{Product Owner \& Team Lead}{Jun 2021 -- Feb 2022}
\begin{itemize}
  \item Construí una plataforma marketplace en el stack MERN (MongoDB, Express, React, Node)
        con servicios GCP (Firebase, Firestore, BigQuery).
  \item Lideré el descubrimiento de cliente y producto, produje mockups en Figma y gestioné
        la entrega con Jira, Confluence y ProjectLibre.
\end{itemize}

%================= EDUCACIÓN ==================================
\section{Educación}
\entry{Universidad Católica del Norte}{Antofagasta, Chile}{Magíster en Ciencias de la Computación}{2019}
\vspace{3pt}
\entry{Universidad Católica del Norte}{Coquimbo, Chile}{Ingeniero Civil en Computación e Informática}{2018}
\vspace{3pt}
\entry{Universidad Católica del Norte}{Coquimbo, Chile}{Licenciado en Ciencias de la Ingeniería}{2016}

%================= HABILIDADES ================================
\section{Habilidades y Certificaciones}
\begin{itemize}
  \item \textbf{AWS / GCP:} Lambda, API Gateway, EventBridge, SQS, SNS, Kinesis, Step Functions,
        DynamoDB, RDS, Aurora, S3, VPC/Endpoints, EC2, Fargate, CloudFormation, Service Catalog,
        Secrets Manager, Certificate Manager, Powertools, CloudWatch, Bedrock, Bedrock Agents,
        IoT Core, Glue, QuickSight, WAF, Route\,53, CloudFront; Firebase, BigQuery,
        Compute Engine, Cloud Run.
  \item \textbf{Lenguajes y Frameworks:} Python, Ruby on Rails, Django, JavaScript, TypeScript,
        Node/Express, React, React Native, Vite, Perl, \LaTeX.
  \item \textbf{Bases de Datos:} MySQL, PostgreSQL, Oracle Database, DynamoDB, MongoDB, Firestore; análisis SQL.
  \item \textbf{Prácticas y Herramientas:} Serverless Framework, GitHub Actions CI/CD, Docker, Make, pytest,
        Grafana, PostHog, MS Clarity, Figma, Jira, OKR, Agile (6+ años), Lean Six Sigma.
  \item \textbf{Certificaciones:} AWS Cloud Practitioner; Refactorización de Software
        (Refactoring Guru); Scrum Fundamentals (SCRUMstudy).
  \item \textbf{Idiomas:} Inglés (IELTS B2), Español (nativo), Portugués (básico), Japonés (básico).
\end{itemize}

\end{document}
